The train and test accuracies are given in \autoref{fig:pdl1_experiment_2_cross_train_test_mw}.
It shows results for $Q_y=0.5$ (top row) and $Q_y=0.6$ (bottom row).
The first column corresponds to $\beta=0$.
Hence, the first column is a variant of \autoref{fig:threshold_train_test_accuracy_mw_0.svg}.
The margin width $\beta$ increases with the column index, the last column corresponding to $\beta=1$.

From this plot, there is no clear improvement in test accuracy for $\beta>0$.
Increasing $\beta$ does appear to dampen the effect of $C$ on train accuracy.
For larger $\beta$, increasing $C$ leads to less improvement in train accuracy.
Hence, the ramp loss model may be less prone to overfitting.
It may also indicate that the ramp loss model is less capable of learning altogether.
Considering the confidence intervals for test accuracy,
no significant improvement in test accuracy can be established.
There may even be a decrease in test accuracy as $\beta$ increases.

Based on the plot, we study $\beta=0.6$ in more detail.
The ruleset frequency table is given in \autoref{tab:threshold_cross_interpretable_mw_0.6_C_1}.
For $Q_y=0.5$, the ruleset found is now consistent among all splits.
The opposite is true for $Q_y=0.6$.
Interestingly, the reverse was true for $\beta=0$, as seen in \autoref{tab:threshold_cross_interpretable_mw_0_C_1}.
It appears therefore that the introduction of the ramp loss,
which lead to the development of the non-additive model, did not lead to more consistent rulesets being found.
Hence, there appears to be no advantage of using the ramp loss model over the threshold model.

\clearpage

\begin{landscape}
    \vspace*{\fill}
    \begin{figure}[h]
        \includesvg[width=\textwidth]{gen/PDL1ex2cross/train_test_mw.svg}
        \caption{Confidence intervals for train and test accuracies for the ramp loss model, $C=1,2,3$, $D=2$, \\ $\beta=0,0.2,0.4,0.6,0.8,1$.}
        \label{fig:pdl1_experiment_2_cross_train_test_mw}
    \end{figure}
    \vspace*{\fill}
\end{landscape}


\begin{table}[h]
    \centering
    \begin{tabular}{llllr}
\toprule
 &  &  &  & \textbf{Frequency} \\
$\mathbf{Q_y}$ & $\mathbf{C}$ & $\mathbf{D}$ & \textbf{TF combination} &  \\
\midrule
\multirow[t]{2}{*}{0.5} & \multirow[t]{2}{*}{1} & 2 & {RELA $\wedge$ STAT2} & 10 \\
\cline{3-5}
 &  & 3 & {FOS $\wedge$ RELA $\wedge$ STAT2} & 10 \\
\cline{1-5} \cline{2-5} \cline{3-5}
\multirow[t]{6}{*}{0.6} & \multirow[t]{6}{*}{1} & \multirow[t]{2}{*}{2} & {RELA $\wedge$ STAT2} & 5 \\
 &  &  & {FOSL1 $\wedge$ RELA} & 5 \\
\cline{3-5}
 &  & \multirow[t]{4}{*}{3} & {FOS $\wedge$ RELA $\wedge$ STAT2} & 4 \\
 &  &  & {FOSL1 $\wedge$ RELA $\wedge$ STAT2} & 3 \\
 &  &  & {FOSL1 $\wedge$ IRF1 $\wedge$ RELA} & 2 \\
 &  &  & {FOS $\wedge$ FOSL1 $\wedge$ RELA} & 1 \\
\cline{1-5} \cline{2-5} \cline{3-5}
\bottomrule
\end{tabular}

    \caption{}
    \label{tab:threshold_cross_interpretable_mw_0.6_C_1}
\end{table}